\documentclass[12pt,a4paper]{article}
\usepackage[utf8]{inputenc}
\usepackage[spanish]{babel}
\usepackage{geometry}
\geometry{a4paper, margin=2.5cm}
\usepackage{graphicx}
\usepackage{booktabs}
\usepackage{tabularx}
\usepackage{hyperref}

\begin{document}

\begin{center}
\vspace*{1cm}
{\LARGE \textbf{Campanil Challenge: Programación Competitiva}} \\
\vspace{0.5cm}
{\large \textbf{Universidad de Concepción - IEEE RAS UdeC}} \\
{\large \textbf{5 de diciembre de 2026}} \\
\end{center}
\vspace{1cm}

\section{Objetivo de la Competencia}
La competencia de Programación Competitiva del Campanil Challenge tiene como objetivo poner a prueba las habilidades de resolución de problemas, pensamiento lógico y trabajo en equipo de los estudiantes, a través del desarrollo de algoritmos en un entorno controlado y con límite de tiempo.

\section{Equipos y Participantes}
2.1. Los equipos estarán conformados por un máximo de \textbf{3 participantes}. \\
2.2. Cada equipo debe designar a un/a Capitán/a, quien será el encargado de la comunicación principal con los organizadores y jueces. \\
2.3. Ningún participante puede pertenecer a más de un equipo en esta misma categoría.

\section{Inscripción}
3.1. Las inscripciones estarán abiertas hasta el \textbf{22 de noviembre de 2026 a las 23:59 hrs}. \\
3.2. La inscripción y la participación son gratuitas. No se requiere pago ni comprobante. \\
3.3. El proceso de inscripción se realizará a través del formulario oficial dispuesto por la organización del evento.

\section{Formato del Concurso}
4.1. \textbf{Duración:} El concurso tendrá una duración oficial de \textbf{4 horas} continuas. \\
4.2. \textbf{Problemas:} Se entregará un conjunto de problemas (entre 6 y 10 problemas algorítmicos) en español y/o inglés. \\
4.3. \textbf{Equipamiento:} Se permitirá el uso de \textbf{una (1) computadora} por equipo. Los equipos podrán llevar su propio teclado y ratón si lo desean, previa revisión de los jueces. \\
4.4. \textbf{Lenguajes permitidos:} Se podrá programar en \textbf{C, C++, Java y Python 3}. Las versiones específicas y entornos de desarrollo serán confirmados a los equipos inscritos antes del evento. \\
4.5. \textbf{Material de apoyo:} Los equipos pueden utilizar material impreso (libros, apuntes, diccionarios). \textbf{No está permitido} el uso de dispositivos electrónicos adicionales (teléfonos, tablets, calculadoras inteligentes, relojes inteligentes) ni memorias USB con código pre-escrito. \\
4.6. \textbf{Acceso a Internet:} Durante la competencia, el acceso a Internet estará estrictamente limitado a la plataforma de evaluación. El uso de motores de búsqueda, foros, herramientas de inteligencia artificial (ej. ChatGPT, Copilot) o comunicación con personas externas resultará en la \textbf{descalificación inmediata} del equipo.

\section{Evaluación y Puntuación}
5.1. El envío de soluciones se realizará mediante una plataforma de evaluación automática (Juez Virtual). \\
5.2. El jurado y la plataforma determinarán si el código enviado es correcto ("Accepted") o si incurre en errores ("Wrong Answer", "Time Limit Exceeded", "Run-Time Error", etc.). \\
5.3. El ganador de la competencia será el equipo que resuelva la mayor cantidad de problemas. \\
5.4. \textbf{Desempate:} En caso de empate en la cantidad de problemas resueltos, se utilizará el \textbf{Tiempo de Penalización} (Penalty Time). El equipo con menor tiempo de penalización será el ganador.
\begin{itemize}
    \item El tiempo de penalización es la suma de los minutos transcurridos desde el inicio de la competencia hasta que se envía una solución aceptada para cada problema resuelto.
    \item Cada envío incorrecto de un problema que finalmente es resuelto suma una penalización adicional de \textbf{20 minutos}.
    \item Los problemas no resueltos no acumulan tiempo de penalización.
\end{itemize}

\section{Reglas de Conducta}
6.1. Se espera que todos los competidores actúen con integridad y respeto hacia los demás participantes, jueces y organizadores. \\
6.2. Cualquier intento de sabotaje a la red, a la plataforma de evaluación o a otros equipos será sancionado con la expulsión del evento. \\
6.3. Las decisiones emitidas por el comité de jueces son \textbf{inapelables}.

\vspace{1cm}
\noindent El equipo organizador se reserva el derecho a modificar las bases y reglas de la competencia de ser necesario, lo cual será informado pertinentemente a todos los equipos inscritos.

\end{document}
